\section{UniLIP-CS2 Benchmark}
\label{sec:benchmark}

We construct a map-aware benchmark for learning the correspondence between
an overhead tactical radar map, a first-person view (FPV), and the camera
pose in Counter-Strike 2 (CS2).  Each example contains a map-specific radar
image, an FPV image captured from the running game, and a five-dimensional
pose consisting of two horizontal coordinates, a vertical coordinate, pitch,
and yaw.  The benchmark supports both directions of this correspondence:
localizing the camera from an FPV/radar pair and synthesizing an FPV image
from a radar map and a specified pose.  It also contains ordered short
trajectories, allowing image generation to be evaluated on temporally
continuous motion rather than only on independently sampled frames.

The construction is a deterministic pipeline after an explicit audit: game
capture produces paired visual and state observations; map resources provide
the radar coordinate system; preprocessing reconciles the two state streams;
quality control fixes the retained source rows and per-map height ranges; and
the final split is formed from complete map--\texttt{file\_num} records.  The
resulting manifests, calibration files, build report, and checksums are the
definition of the benchmark instances used in the experiments.
Figure~\ref{fig:cs2-benchmark-pipeline} summarizes this pipeline.

\begin{figure*}[t]
    \centering
    \fbox{\parbox[c][3.35cm][c]{0.94\linewidth}{\centering
    \textit{Placeholder for the benchmark construction figure.}\\[2pt]
    The final figure should show, from left to right, the controlled CS2
    setup and the three capture streams (FPV screenshot, GSI state, and
    process-memory state), radar extraction and map-coordinate conversion,
    Mem/GSI consistency filtering, missing-image removal, and audit,
    full-corpus height calibration, \texttt{file\_num}-record pool assignment,
    and the final localization and conditional-generation tasks.  It should
    distinguish individual-frame files from ordered 64-frame continuous
    clips.}}
    \caption{Placeholder.  The final figure should visualize the complete
    data path and identify the artifact produced at each stage, including
    the point at which record-level isolation is enforced.}
    \label{fig:cs2-benchmark-pipeline}
\end{figure*}

\subsection{Data Acquisition and Canonical Representation}

\paragraph{Controlled CS2 capture.}
The data are collected from a CS2 client in a controlled local
practice/spectator setup.  We manually launch the game from Steam with
\texttt{-console -insecure} and use a windowed 4:3, 1024$\times$768 display.
Console commands hide the selected HUD overlays,
weapon/view model, and crosshair so that the recorded FPV image exposes the
game scene without those interface elements.  The recorder observes the
selected player from the spectator view while the game is running.

Under the Game State Integration (GSI) configuration, the CS2 client sends
JSON updates to a local HTTP endpoint at \texttt{localhost:3000}, with
\texttt{player\_position} enabled together with the relevant map and player
state.  GSI provides the
player position and game metadata but does not expose the complete camera
viewpoint.  The recorder therefore reads the running game process to obtain
the position and the pitch/yaw viewpoint fields.  In each capture-loop
iteration it associates a whole-window screenshot with the most recent GSI
payload and a process-memory readout.  These observations are paired by the
recorder loop; the implementation does not provide a shared hardware
timestamp or claim timestamp-atomic synchronization.

Raw observations are stored in \texttt{record\_N.npy} batches.  Each entry
contains the FPV image path and the memory and GSI dictionaries, while the
images are retained as individual JPEG files.  The resulting identity,
\texttt{file\_numN\_frame\_M}, preserves both the record-batch identifier and
the frame identifier.  Here, a record denotes one \texttt{record\_N.npy}
batch identified by \texttt{file\_num}; it is not a physical capture-session
identifier.  The complete map--\texttt{file\_num} group is carried through
preprocessing and later serves as the atomic split unit.

\paragraph{Radar maps and coordinate transform.}
The radar image for each map is obtained from the CS2 resource package; the
data-collection documentation describes its extraction with Source 2 Viewer.
The corresponding map-overview metadata supplies an offset and scale that
align game-world positions with the radar image.  Let
$(x_w,y_w,z_w)$ denote the camera position, where $z_w$ is the camera eye
height, and let $(o_x,o_y,s)$ denote the map-specific offset and scale.  We
convert the position to the canonical radar coordinates
\begin{equation}
    x_r = \frac{x_w-o_x}{s}, \qquad
    y_r = \frac{o_y-y_w}{s}, \qquad
    z_r = \frac{z_w}{s}.
\end{equation}
The reversal in the second equation accounts for image coordinates.  The
implementation casts the transformed coordinates to integer radar units;
$x_r$ and $y_r$ index a 1024$\times$1024 radar image, while $z_r$ is retained
as the vertical coordinate.  The upper- and lower-floor radar layers are
blended for Nuke, Train, and Vertigo, as specified by the map assets.  Radar
paths are explicit per-map entries, so a missing asset cannot be silently
replaced by an asset from another map.

\paragraph{State reconciliation and row schema.}
The preprocessing code independently transforms the memory position and the
GSI position into radar coordinates.  A frame is removed when the Euclidean
distance between the two horizontal positions exceeds 3 radar pixels.  A
frame whose referenced FPV image is missing is also removed.  For every
retained frame, the image is assigned a stable name and the metadata are
written to the map's \texttt{positions.json}.  A canonical row is
\begin{equation}
    (x_r,y_r,z_r,\theta_h,\theta_v,\mathrm{map},\mathrm{file\_frame}),
\end{equation}
where $\theta_h$ is yaw and $\theta_v$ is pitch, both stored in radians.
Thus, a row identifies one FPV JPEG, one map-specific radar image, and one
five-dimensional pose without discarding the \texttt{file\_num} identity.

\subsection{Quality Control and Benchmark Construction}
\label{subsec:data_filter}

\paragraph{Integrity audit.}
The construction first validates frame identities, duplicate identities,
finite pose values, map labels, source images, and radar assets.  It then
reports coordinate candidates using out-of-window coordinates and large
adjacent-frame changes in position or angle.  These conditions are review
triggers rather than unconditional rejection rules: multi-level maps and
camera motion can produce large valid changes.  A reviewed decision file
records the default action and explicit exceptions, making the retained
source corpus reproducible.

The production audit contains 10,345 coordinate candidates and no
integrity-invalid rows.  The approved decisions retain 9,705 candidate rows
under the declared default action and exclude 18,756 rows, including 19
complete map--\texttt{file\_num} records.  After these decisions, the source
corpus contains 3,087,392 retained rows in 3,430
map--\texttt{file\_num} records across the
14 maps.  Of these rows, 2,280,253 belong to the ten seen maps and 807,139 to
the four held-out maps.

\paragraph{Frozen height calibration.}
After the audit and before any record is assigned to a split, the builder
computes an exact minimum and maximum of $z_r$ independently for every map
over the entire approved source corpus.  The task representation uses
\begin{equation}
    z' = \frac{z_r-z_{\min,m}}{z_{\max,m}-z_{\min,m}},
\end{equation}
where $m$ is the map.  The bounds are frozen, the normalized range is
$[0,1]$, no clipping is applied, and a missing or degenerate range raises an
error instead of invoking a split-derived fallback.  Because the calibration
is computed before splitting, the same map-specific convention is used for
training, validation, support, query, and continuous rows.  In particular,
the supplied calibration for a held-out map is part of the benchmark input;
held-out-map results should therefore be interpreted as zero-shot model
transfer and few-shot adaptation under fixed map-level calibration, not as
calibration-free transfer.

\paragraph{Map roles.}
The 14 map sources are divided into a ten-map seen set and a four-map
held-out set.  Seen-10 contains
\texttt{cs\_agency}, \texttt{cs\_italy}, \texttt{de\_ancient},
\texttt{de\_anubis}, \texttt{de\_dust2}, \texttt{de\_inferno},
\texttt{de\_mirage}, \texttt{de\_nuke}, \texttt{de\_overpass}, and
\texttt{de\_train}.  CrossMap-4 contains \texttt{cs\_office},
\texttt{de\_golden}, \texttt{de\_palacio}, and \texttt{de\_vertigo}.
Each manifest row retains the map identity, while the manifest provides one
explicit radar asset path for each map.

\paragraph{Record-level split construction.}
The complete map--\texttt{file\_num} group is the atomic split unit.  All
frames from one such record are assigned to one pool before any frame
sampling, preventing rows from the same record from entering different
pools.  Records are ranked with SHA-256 using the
fixed global seed 20260827 and allocated with deterministic largest-remainder
rounding.  The record-pool ratios are 60\%/10\%/20\%/10\% for
Seen-10 train/validation/discrete-test/continuous and 20\%/60\%/20\% for
CrossMap-4 support/query/continuous.

For Seen-10 test selection, near-pose filtering is first applied against the
selected train and validation poses; for CrossMap-4 query selection, it is
first applied against the union of all five support draws.  The tolerances are
8 radar units in horizontal-plane Euclidean distance, 8 units in the vertical
coordinate, and 5 degrees for each angle, with circular distance for yaw.
Eligible frames for every discrete split are then ordered within each record,
thinned to a minimum five-frame gap, and selected by deterministic coverage
sampling over $4\times4\times3\times8\times3$ quantile bins in
$(x_r,y_r,z_r,\theta_h,\theta_v)$.  A quota failure stops the construction;
the builder does not silently reduce a requested split.

Continuous clips are selected only from the continuous record pool.  Each
map contributes 20 clips, each clip contains exactly 64 ordered frames, and
successive frame identifiers differ by one or two.  No frame is reused within
the continuous split, and at most one clip is selected from a record.  Clip
identifiers and frame order are retained in the manifest, so continuous
evaluation uses the recorded sequence rather than reconstructing tracks from
unordered filenames.

\subsection{Final Splits and Dataset Statistics}
\label{subsec:benchmark-splits}

For every Seen-10 map, the final discrete split contains 5,000 training
frames, 500 validation frames, and 2,000 discrete test frames.  Across the
ten maps this gives 50,000 training frames, 5,000 validation frames, and
20,000 discrete test frames.  The continuous split contributes 20 clips of
64 frames per map, or 1,280 frames per map and 12,800 frames in total.

For every CrossMap-4 map, each support draw contains 100 frames, the fixed
query contains 2,000 frames, and the continuous split contains 20 clips of
64 frames.  Five deterministic support draws (seeds 0--4) are provided.  A
single support draw therefore contains 400 frames over the four maps, while
the fixed query contains 8,000 frames and the continuous set contains 5,120
frames.  Support draws may overlap one another; the query and continuous
sets are identical across support seeds.  The union of all support records is
disjoint from the query and continuous records, which are also disjoint from
one another.  Counting all five support draws separately gives 102,920
manifest row references; overlap among support draws yields 102,780 unique
selected FPV images.

The pre-audit source snapshot contains 3,106,148 preprocessed rows and 3,449
source records; after quality control, 3,087,392 rows and 3,430 records
remain.  Seen-10 contributes 2,280,253 retained rows and CrossMap-4
contributes 807,139.  The map-level source and retained counts are reported
together with the map-specific $z$ extrema because maps differ in the number
of records and vertical span.
Tables~\ref{tab:benchmark-statistics} and~\ref{tab:benchmark-splits}
summarize the map-level corpus and the final task splits, respectively.

\begin{table*}[t]
    \centering
    \fbox{\parbox[c][4.05cm][c]{0.94\linewidth}{\centering
    \textit{Placeholder for the map-level dataset statistics table.}\\[2pt]
    The final table should contain one row for each of the 14 maps and report
    its Seen-10/CrossMap-4 role, preprocessed rows, approved-retained rows,
    source and retained \texttt{file\_num} records, frozen $z_{\min}$ and
    $z_{\max}$, and the exact radar asset (including blended assets).  An
    aggregate row should report 3,106,148 source rows, 3,087,392 retained
    rows, 3,449 source records, and 3,430 retained records.}}
    \caption{Placeholder.  The final table should make map imbalance,
    retained-corpus size, vertical calibration ranges, and radar-asset
    choices auditable.}
    \label{tab:benchmark-statistics}
\end{table*}

\begin{table*}[t]
    \centering
    \fbox{\parbox[c][3.75cm][c]{0.94\linewidth}{\centering
    \textit{Placeholder for the final split and task table.}\\[2pt]
    The final table should report, for Seen-10 and CrossMap-4, the map role,
    record-pool ratio, number of records, discrete rows or support/query
    rows, number of continuous clips, frames per clip, support seeds, and
    the exact task/evaluation split.  It should explicitly state the
    5,000/500/2,000 Seen-10 quotas, the 100-shot-per-map CrossMap quota, the
    five support seeds, and the fixed 64-frame continuous clips.}}
    \caption{Placeholder.  The final table should connect each split to its
    intended task and document all quotas and disjointness constraints.}
    \label{tab:benchmark-splits}
\end{table*}

\subsection{Tasks and Evaluation Protocol}
\label{subsec:benchmark-tasks}

\paragraph{Task inputs and outputs.}
The benchmark defines two frame-level directions and one sequence-level
evaluation setting.  For localization, the inputs are the FPV image, the
corresponding overhead radar image, and the map identity; the output is the
five-dimensional pose $(x,y,z,\mathrm{pitch},\mathrm{yaw})$.  For discrete
conditional generation, the inputs are the radar image, map identity, and a
specified pose; the output is the corresponding FPV image.  Continuous
conditional generation uses the same conditioning variables for each frame
but evaluates the generated images in the exact ordered 64-frame clips.
The benchmark does not prescribe whether a model trains the two branches
jointly or trains either branch alone; the selected manifest and split must
always be reported.

The on-disk pose fields are $(x_r,y_r,z_r,\theta_h,\theta_v)$, with yaw
stored as $\theta_h$ and pitch stored as $\theta_v$.  The task serialization
uses the order $(x,y,z,\mathrm{pitch},\mathrm{yaw})$ and the normalized tuple
\begin{equation}
    \left(\frac{x_r}{1024},\frac{y_r}{1024},z',
    \frac{\theta_v}{2\pi},\frac{\theta_h}{2\pi}\right).
\end{equation}
The prompt convention defines yaw zero as east with clockwise increase,
pitch zero as looking straight down and pitch $\pi$ as looking straight up.
Keeping the storage names and the task-token order explicit avoids swapping
pitch and yaw during training or evaluation.

\paragraph{Seen-10 evaluation.}
The base model is trained only on Seen-10 train records.  Seen-10 validation
records are reserved for checkpoint selection and hyperparameter decisions;
they are not used to report the final test result.  Localization and
discrete generation are evaluated on \texttt{seen\_discrete\_test}, while
continuous generation is evaluated on \texttt{seen\_continuous}.  The
individual test rows and continuous clips are record-disjoint from training
and validation under the split construction above.

\paragraph{CrossMap-4 transfer.}
CrossMap-4 maps are excluded from base-model optimization.  Zero-shot
transfer evaluates localization and discrete generation on the fixed
\texttt{crossmap\_query\_test} rows and evaluates continuous generation on
the fixed \texttt{crossmap\_continuous} clips.  In the 100-shot setting, the
same Seen-10 checkpoint is used to initialize each of the five support-seed
runs.  Adaptation uses exactly 100 labeled frames per held-out map, or 400
frames in total, and is then evaluated on the same fixed query and
continuous sets.  An adapted run never initializes from another support
seed, and query or continuous rows cannot be used for checkpoint selection,
early stopping, or adaptation-budget tuning.

All results are reported per map and as an equal-map macro average.  The
five support draws are summarized by their mean and a two-sided 95\%
Student-$t$ confidence interval; variation from independent model-training
seeds is a separate uncertainty source.  Generation coverage uses the
manifest-selected rows as its denominator.  For sequence metrics, the
64-frame clip boundaries are preserved; in particular, Fr\'echet Video
Distance (FVD) uses 16-frame windows with stride 16 within each clip and never
forms a window across two clips.

\subsection{Reproducibility and Scope}
\label{subsec:benchmark-scope}

The benchmark bundle contains the per-map split files, aggregate
manifests, map-specific height calibration, audit/build metadata, and
checksums for the generated artifacts and selected FPV images.  The manifest
also records the map lists, global and support seeds, record pools, exact
quotas, source fingerprints, radar paths, calibration fingerprints, and
selected image identities.  This makes changes to the bound source rows,
radar assets, calibration, or selected images detectable without implying
that unselected files are part of the benchmark.

The strongest trajectory-isolation claim supported by the data is
\texttt{file\_num}-record disjointness: the preprocessed metadata do not
contain a capture-session identifier, so record disjointness should not be
described as physical-session disjointness.  Continuous clips are
record-disjoint but are not guaranteed to be pose-disjoint from every
training frame.  CrossMap-4 therefore measures transfer to held-out map
records under the supplied radar assets, map identity, and frozen map-level
height calibration.  It does not by itself establish domain generalization
without target-map metadata.  These boundaries, together with the fixed
record pools and manifest-selected denominators, define the claims that can
be made from the benchmark.
